\section{Model}
\label{sec:model}

\begin{figure}[!t]
  \centering
  \includegraphics[width=0.92\linewidth]{figures/architecture.pdf}
  \caption{AmphionASR architecture. The audio branch (left) maps a
    16\,kHz waveform through a Whisper-style log-mel front-end, a
    convolutional sub-sampler, a 24-layer Whisper-style Transformer
    encoder and a lightweight modality projector into the LLM's
    2560-dimensional hidden space. The text branch (right) tokenises
    the chat template and replaces four special boundary tokens with
    learnable prompt vectors. The two streams meet at a merge step that
    substitutes the \texttt{<speech>} placeholder with the audio
    embedding sequence, after which Qwen3-4B-Instruct performs
    next-token prediction on the assistant span only.}
  \label{fig:architecture}
\end{figure}

\begin{figure}[!t]
  \centering
  \includegraphics[width=\linewidth]{figures/retrieve.pdf}
  \caption{Retrieve-then-bias paradigm. Left: AmphionASR is invoked
    twice -- a no-hotword coarse pass $g(x)$ produces a coarse
    transcript that a parameter-free text retriever maps to the
    top-$K$ subset of a hotword pool $\mathcal{H}$; a
    discriminative-gap check refines the list; and a with-hotword
    fine pass $f(x, \mathcal{H}')$ produces the final transcript.
    Right (dashed): the audio-space alternative
    of~\cite{min2024speechrag}, not adopted by AmphionASR.}
  \label{fig:retrieve}
\end{figure}

\subsection{Overview}

AmphionASR follows the now-standard SpeechLLM template
\cite{ma2024slamasr, chu2023qwenaudio, tang2024salmonn} -- an audio
encoder, a small modality projector and a frozen-then-fine-tuned
LLM -- but is specialised in three respects that we describe in this
section: (i)~a Whisper-style audio encoder initialised from the
Qwen3-ASR audio tower; (ii)~a lightweight two-layer projector that
uses the SwooshR activation \cite{yao2024zipformer}; and (iii)~four
learnable \emph{modality-boundary prompts} that mark the audio--text
boundary inside the LLM's input-embedding stream. The full
computational graph is shown in \Cref{fig:architecture} and the
hotword retrieval pipeline in \Cref{fig:retrieve}.

Component-level parameter counts are summarised in
Table~\ref{tab:params}. The total of $\approx\!4.74$\,B includes the
tied output projection; counting distinct parameters the model is
$\approx\!4.3$\,B.

\begin{table}[t]
  \caption{Per-component parameter counts.}
  \label{tab:params}
  \centering
  \small
  \begin{tabular}{lrl}
    \toprule
    Component & Parameters & Role \\
    \midrule
    Audio encoder              & $314.33$\,M     & Whisper-style 24-layer Transformer \\
    Modality projector         & $9.18$\,M       & Two-layer MLP with SwooshR activation \\
    Language model             & $4{,}411.42$\,M & Qwen3-4B-Instruct (Qwen3-4B-Instruct-2507) \\
    Boundary prompts           & $0.01$\,M       & $4\!\times\!2560$ learnable vectors \\
    \midrule
    Total                      & $\approx 4.735$\,B & $\approx 4.3$\,B distinct after input--output embedding tying \\
    \bottomrule
  \end{tabular}
\end{table}

\subsection{Audio Encoder}
\label{subsec:audio-encoder}

The audio encoder is initialised from the AuT Encoder of Qwen3-ASR
\cite{qwen2025qwen3asr}, a 24-layer Whisper-style pre-LN Transformer.
It accepts 16\,kHz mono audio, processes it through a log-mel front-end
(128 mel bins, 10\,ms hop) and a convolutional sub-sampler, and
produces 1024-dimensional frame representations at ${\approx}13$\,Hz.
A chunked-inference window of 800 frames extends the effective input
duration beyond the standard Whisper 30\,s ceiling, enabling long-form
transcription without chunk stitching.

\subsection{Modality Projector}
\label{subsec:projector}

The modality projector is a lightweight two-layer MLP that adapts the
encoder output to the LLM's input requirements: it maps the
1024-dimensional AuT Encoder output frames into the LLM's
2560-dimensional semantic embedding space, bridging the two modalities
without any temporal sub-sampling. It contributes about $9.2$\,M
parameters and is the only fully-trainable module in Stage~1
(Section~\ref{sec:training}).

\subsection{Language Model}
\label{subsec:llm}

The decoder is Qwen3-4B-Instruct (the Qwen3-4B-Instruct-2507 release
of the Qwen3 family \cite{qwen2025qwen3}), a 36-layer decoder-only
Transformer with 2560-dimensional hidden states and a 256\,K native
context length. Its 256\,K context carries over to AmphionASR, enabling
long-form transcription without chunk stitching. The model uses
input--output embedding tying; the component count in
Table~\ref{tab:params} therefore over-counts by one embedding tensor,
hence ``$\approx\!4.3$\,B distinct'' alongside
``$\approx\!4.74$\,B total''. The released checkpoint folds the LoRA
adapters trained in Stages~2 and~3 (Section~\ref{sec:training}) back
into the base weights, so the published weight files contain no
separate LoRA, adapter or quantisation tensors.

\subsection{Special Tokens and Chat Template}
\label{subsec:tokens}

In addition to the Qwen3 vocabulary, AmphionASR defines the five
special tokens listed in Table~\ref{tab:tokens}. The audio segments
are wrapped as \verb|<start_speech><speech><end_speech>|, where
\verb|<speech>| is the placeholder that the merge step replaces
with the projected audio embedding sequence.

\begin{table}[t]
  \caption{Special tokens defined on top of the Qwen3 vocabulary.}
  \label{tab:tokens}
  \centering
  \small
  \begin{tabular}{lcl}
    \toprule
    Token            & Id      & Role \\
    \midrule
    \texttt{<speech>}        & 151669 & audio placeholder, replaced by audio embeddings \\
    \texttt{<start\_text>}    & 151670 & text-segment start; embedding overridden \\
    \texttt{<end\_text>}      & 151671 & text-segment end \\
    \texttt{<start\_speech>}  & 151672 & speech-segment start \\
    \texttt{<end\_speech>}    & 151673 & speech-segment end \\
    \bottomrule
  \end{tabular}
\end{table}

Each task is exposed to the LLM as a short instruction line followed
by the audio placeholder (Table~\ref{tab:tasks}). The language hint
and the hotword list are written as additional lines after the
instruction; the target-speaker variant prepends an enrollment audio
block before the main audio. The same prompt template is used at
training time and at inference time; a mismatch between the two
causes the model to either repeat hotwords verbatim or skip the
language hint, with measurable WER damage.

\begin{table}[t]
  \caption{Task prompt templates used by AmphionASR. \texttt{\{speech\}}
    expands to \texttt{<start\_speech><speech><end\_speech>} via the
    chat template; \texttt{\{enroll\}} expands similarly to the
    enrollment utterance. Bracketed lines are optional.}
  \label{tab:tasks}
  \centering
  \small
  \setlength{\tabcolsep}{4pt}
  \begin{tabular}{lp{9.6cm}}
    \toprule
    Task & Prompt template \\
    \midrule
    Plain ASR        & \texttt{Transcribe the following audio.\{speech\}} \\
    ASR (English)    & \texttt{Transcribe the following audio.\textbackslash nLanguage: English\{speech\}} \\
    ASR (Mandarin)   & \texttt{Transcribe the following audio.\textbackslash nLanguage: Chinese\{speech\}} \\
    Hotword ASR      & \texttt{Transcribe the following audio.[\textbackslash nLanguage: ...]\textbackslash nHotwords:\,hw1,hw2\{speech\}} \\
    TS-ASR           & \texttt{Given the speaker's voice:\{enroll\}\textbackslash nTranscribe what this speaker says in the following audio.[\textbackslash nLanguage: ...][\textbackslash nHotwords: ...]\{speech\}} \\
    \bottomrule
  \end{tabular}
\end{table}

\paragraph{Target-speaker conditioning.}
TS-ASR is implemented entirely in the prompt and the chat template:
the enrollment utterance is supplied as an additional audio content
item that the chat template renders as a separate
\verb|<start_speech><speech><end_speech>| triplet, and the main
instruction switches from ``\textit{Transcribe the following
audio.}'' to ``\textit{Transcribe what this speaker says in the
following audio.}''. The model itself is unchanged -- both audio
segments share the audio encoder and projector path and reach the
LLM as two ordered audio-embedding spans (enrollment first, mixed
audio second). Anti-hallucination negative samples
(Section~\ref{subsec:tsasr-synth}) inherit the same template with an
empty assistant response.

\paragraph{Hotword conditioning.}
Hotwords are inserted into the prompt as a single comma-separated line.
At training time the line is included with probability $0.8$; at
evaluation time the eval client pads each utterance to a target $K$
hotwords by mixing the ground-truth list with random distractors drawn
from the dataset's own hotword pool (Section~\ref{subsec:hotword-pool}).
No special token marks the hotword list, so the LLM must learn the
\verb|Hotwords:| literal as part of the prompt grammar.

\FloatBarrier
